\chapter{Внутреннее происхождение ковариации межсекторных стрелок}
\label{ch:v8-cross-arrow-covariance-origin-gate}

\section{Что может дать уже известный гессиан}

Предыдущий гейт оставил ненулевую ковариацию как необходимое условие
работы Краус-моста. Проверим три правила, не добавляющие новых стрелок:
\begin{equation}
 C_{\rm cl}=a^{-1}H_\eta^{-1},\qquad
 C_{\rm q}=\frac{\hbar}{2\sqrt Z}H_\eta^{-1/2},\qquad
 C_\tau=e^{-\tau H_\eta},
 \label{eq:v8-three-internal-covariance-candidates}
\end{equation}
где
\begin{equation}
 H_\eta=H_E+2\eta H_{\rm link}.
\end{equation}
Первая формула является конечномерной гауссовской мерой, вторая ---
основным состоянием гармонического приближения при скалярной кинетической
метрике $Z$, третья --- полным корреляционным ядром квадратичного
генератора.

\section{Новая положительная структура}

На двенадцати межсекторных стрелок направлениях связывающий гессиан распадается в шесть
тождественных вещественных пар $(QLYR,XLdR)$:
\begin{equation}
 H_{\rm link}^{qX}
 =I_6\otimes
 \begin{pmatrix}
 0.92492997&-0.44926187\\
 -0.44926187&0.58178422
 \end{pmatrix}.
 \label{eq:v8-linking-repeated-cross-pair}
\end{equation}
Его собственные значения равны $0.27244808$ и $1.23426610$, а связь с
остальными пятнадцатью вещественными направлениями точно отсутствует на
проверенном срезе.

Следовательно, при любом $\eta>0$ все спектральные функции $H_\eta$
выбирают одну и ту же главную ось смешивания двух типов стрелок. Её угол в
принятом базисе равен
\begin{equation}
 \theta_{qX}=0.96780105\ldots\ {\rm rad}
 =55.45092\ldots^\circ.
 \label{eq:v8-cross-covariance-axis}
\end{equation}
Это не угол CKM: он действует в пространстве типов межсекторных стрелок, а
не между тремя поколениями. Тем не менее это первая ненулевая
корреляционная ось, выбранная уже существующей полярной геометрией.

\section{Почему полной ковариации всё ещё нет}

Ось не определяет эллипс ковариации. При $\eta=0$ обе компоненты полностью
вырождены. При изменении свободного веса $\eta$ отношение главных дисперсий
меняется, хотя собственные направления сохраняются. Поэтому нормированная
форма ковариации не единственна внутри качественного класса Тома VII.

Кроме того, каждое из правил \eqref{eq:v8-three-internal-covariance-candidates}
сохраняет свой внешний масштаб:
\begin{itemize}
 \item классическая мера зависит от общего коэффициента действия $a$;
 \item гармонический вакуум зависит от $\hbar/\sqrt Z$ и требует
 кинетической метрики;
 \item тепловое ядро зависит от корреляционного времени $\tau$.
\end{itemize}
Машинные сканы меняют эти параметры на два порядка и получают
соответствующее изменение Краус-скорости. Нормировка состояния не удаляет
зависимость формы от $\eta$.

\begin{theorem}[Полярная ось без полной меры флуктуаций]
Родитель Тома VII содержит больше информации, чем независимые веса двух
Краус-семейств: относительный связывающая кривизна канонически организует их в шесть
одинаковых пар и при $\eta>0$ фиксирует общую главную ось. Однако ни
анизотропия ковариации, ни её общий масштаб, ни физическая Краус-скорость не
выводятся, поскольку зависят соответственно от $\eta$ и от выбора
классической, квантовой или тепловой меры.
\end{theorem}

\begin{nogocriterion}[Запрет выбора удобной гауссовской меры]
Нельзя положить $a=1$, $\hbar/\sqrt Z=1$ или $\tau=1$ и объявить полученную
дисперсию предсказанием. Эти соглашения дают воспроизводимый представитель,
но не устраняют незамкнутые нормировки Тома VII.
\end{nogocriterion}

\begin{protocol}
\begin{enumerate}
 \item межсекторных стрелок подпространство: \textbf{канонически замкнуто};
 \item одинаковых пар $(QLYR,XLdR)$: \textbf{6};
 \item связь с остальными направлениями: \textbf{ноль};
 \item общая главная ось при $\eta>0$: \textbf{выведена};
 \item её физический статус CKM/PMNS: \textbf{нет};
 \item сила анизотропии: \textbf{зависит от свободного $\eta$};
 \item классический масштаб меры: \textbf{не выведен};
 \item квантовая кинетическая нормировка: \textbf{не выведена};
 \item корреляционное время: \textbf{не выведено};
 \item единственная Краус-скорость: \textbf{не получена};
 \item следующий тест: \textbf{минимальная ковариантная Стайнспринг-дилатация}.
\end{enumerate}
\end{protocol}

Машинный сертификат сохранён в
\path{s2t_v8_cross_arrow_covariance_origin_gate_results.json}.

\begin{thebibliography}{9}
\bibitem{v8-covariance-lewin}
M.~Lewin, P.~T.~Nam, N.~Rougerie,
\emph{Classical field theory limit of many-body quantum Gibbs states in 2D
and 3D}, Inventiones Mathematicae \textbf{224} (2021), 315--444;
arXiv:1810.08370.
\bibitem{v8-covariance-alvarez}
E.~Alvarez,
\emph{Covariant techniques in Quantum Field Theory},
arXiv:2203.11292.
\end{thebibliography}